% Weitere Pakete können problemlos hinzugefügt werden.

% Neue deutsche Rechtschreibung
\usepackage[ngerman]{babel}

\usepackage{csquotes}

% Umbrüche in Zitaten erlauben
\usepackage{cite}

% Zitierungen a la 
% Ziegler (2010) durch citet{ziegler:2010}
% --> t wie textual
% bzw. (Ziegler, 2010) durch \citep{ziegler:2010} 
% --> p wie paranthesis
\usepackage{apacite}
\usepackage{natbib}

% Einbinden von Grafiken
\usepackage{graphicx}

% Komfortableres Tabellen-Paket
\usepackage{tabularx}

% Schriftfamilie Times
\usepackage{palatino}

%Schriftfarbe
\usepackage{xcolor}

% Ermöglicht mehrere Bilder in einer floatenden Figure-Umgebung
\usepackage{subfig}

% Erweiterte Mathe-Zeichen
\usepackage{amsmath,amsmath,amssymb,amsfonts}

% Erlauben von urls
\usepackage{url}

% Einbinden von Listings
\usepackage{listings}

% Klickbare Links innerhalb des PDFs
\usepackage[
	bookmarks=true,         % show bookmarks bar?
    unicode=false,          % non-Latin characters in Acrobat’s bookmarks
    pdftoolbar=true,        % show Acrobat’s toolbar?
    pdfmenubar=true,        % show Acrobat’s menu?
    pdffitwindow=false,     % window fit to page when opened
    pdfstartview={Fit},    % fits the width of the page to the window
    pdftitle={Proposal},    % title
    pdfauthor={\autor},     % author
    pdfsubject={\titel},   % subject of the document
    hidelinks=true,			  % do not show colored links
    pdfdisplaydoctitle=true,	  % show document title instead of filename in titlebar
    pdfpagelayout=TwoColumnRight	% show two pages by default
]{hyperref}

% Anpassung des Abstracts
\usepackage{xpatch}
\xpretocmd{\abstract}{\begin{addmargin}{1.5cm}}{}{}
\xpretocmd{\endabstract}{\end{addmargin}}{}{}
\xapptocmd{\abstract}{\noindent}{}{}